% NOETHER PAPER 2 — KOREAN TRANSLATION-PRODUCER DRAFT — UNCHECKED
% Unit: T04-U17; current-authority whole lines 741--754
% Source slice: 941 LF UTF-8 bytes; SHA-256 61CEA035D29C93A45B9953F0101E9069852785C82662DC8F7718CA2B4DFFF77D
% Pointer: NOETH-DE-AUTH-v004-20260804; SHA-256 A1C62FDACAA34DFC1B806DC18258F2E732539F7AE9D85AA4BD9E1067B8749D9F
% This is producer translation only: no source/Korean/formula review, compilation, or rendering.

전이 원리에 따라 형식 $s_x^m t_x^n$의 $\lambda$번째 편극 $[s_x^m t_x^n]_{y^\lambda}$ ($\lambda\le n$)에는 다음 전개가 성립한다\srcfn{*)}{Gordan, Die Resultante binärer Formen (Kap. I, § 5). Rend. Circ. Mat. di Palermo XXII. 1906.}. $\lambda>n$인 경우에는 $[s_y^m t_y^n]_{x^{m+n-\lambda}}$의 전개를 통하여 첫째 경우로 돌아간다.
\[
\bigl[s_x^m t_x^n\bigr]_{y^\lambda}
=\sum_r (-1)^r\,
\frac{\binom{m}{r}\binom{\lambda}{r}}{\binom{m+n}{r}}\cdot
(st\widehat{xy})^r s_x^{m-r}t_x^{n-\lambda}t_y^{\lambda-r}.
\]
이에 대응하여 $u_\sigma^\mu u_\tau^\nu$의 $\varkappa$번째 편극 $[u_\sigma^\mu u_\tau^\nu]^{v^\varkappa}$에는 다음이 성립한다.
\[
\bigl[u_\sigma^\mu u_\tau^\nu\bigr]^{v^\varkappa}
=\sum_\varrho (-1)^\varrho\,
\frac{\binom{\mu}{\varrho}\binom{\varkappa}{\varrho}}{\binom{\mu+\nu}{\varrho}}\,
(\sigma\tau\widehat{uv})^\varrho u_\sigma^{\mu-\varrho}u_\tau^{\nu-\varkappa}v_\tau^{\varkappa-\varrho}.
\]

